\subsection{Composition product and connecting homomorphisms of extension modules}

PROPOSITION 5. — Let

\[
0 \to {\bf M} ^ {\prime} \stackrel {f} {\to} {\bf M} \stackrel {g} {\to} {\bf M} ^ {\prime \prime} \to 0\tag{8}
\]

be an exact sequence of left A-modules, \(\theta\in\mathrm{Ext}_{\mathbf{A}}^{1}\left(\mathbf{M}^{\prime\prime},\mathbf{M}^{\prime}\right)\) the associated class, N a left A-module, n an integer.

\begin{enumerate}
\def\labelenumi{\alph{enumi})}
\item
  The connecting homomorphism \(\delta^{n}(\mathbf{N},\mathcal{E}):\operatorname{Ext}_{\mathbf{A}}^{n}(\mathbf{N},\mathbf{M}^{\prime \prime})\to \operatorname{Ext}_{\mathbf{A}}^{n + 1}(\mathbf{N},\mathbf{M}^{\prime})\) is the composition product \(\alpha \mapsto \theta \circ \alpha\) by \(\theta\) .
\item
  The connecting homomorphism \(\delta^{n}(\mathcal{E},\mathbf{N}):\operatorname{Ext}_{\mathbf{A}}^{n}(\mathbf{M}^{\prime},\mathbf{N})\to\operatorname{Ext}_{\mathbf{A}}^{n+1}(\mathbf{M}^{\prime\prime},\mathbf{N})\) is the composition product \(\alpha\mapsto(-1)^{n+1}\alpha\circ\theta\) by \((-1)^{n+1}\theta\) .
\item
  Consider a commutative diagram
\end{enumerate}

\[
\begin{array}{l} 0 \xrightarrow {} \mathrm{M} ^ {\prime} \xrightarrow {f} \mathrm{M} \xrightarrow {g} \mathrm{M} ^ {\prime \prime} \xrightarrow {} 0 \\ 0 \xrightarrow {} \mathrm{M} ^ {\prime} \xrightarrow {e _ {\mathrm{M} ^ {\prime}}} \mathrm{I} ^ {0} (\mathrm{M} ^ {\prime}) \xrightarrow {\delta^ {0}} \mathrm{I} ^ {1} (\mathrm{M} ^ {\prime}). \end{array}
\]

By definition, \(\theta\) is the class of \(-v^{1}\in\mathrm{Homgr}_{\mathbf{A}}^{1}(\mathbf{M}^{\prime\prime},\mathbf{I}(\mathbf{M}^{\prime}))\) . Let moreover \(\alpha\in\mathrm{Ext}_{\mathbf{A}}^{n}(\mathbf{N},\mathbf{M}^{\prime\prime})\) , represented by an element \(a\) of \(\mathrm{Homgr}_{\mathbf{A}}^{n}(\mathbf{L}(\mathbf{N}),\mathbf{M}^{\prime\prime})\) . By construction, \(\delta^{n}(\alpha)\) is obtained as follows: one lifts \(a^{n}\in\mathrm{Hom}_{\mathbf{A}}(\mathbf{L}_{n}(\mathbf{N}),\mathbf{M}^{\prime\prime})\) to

\[
b \in \operatorname{Hom} _ {\mathbf {A}} \left(\mathrm{L} _ {n} (\mathrm{N}), \mathrm{M}\right),
\]

and \(\delta^n (\alpha)\) is the class of \(e_{\mathbf{M}'}\circ c\) where \(c\in \mathrm{Hom}_{\mathbf{A}}\left(\mathbf{L}_{n + 1}(\mathbf{N}),\mathbf{M}'\right)\) is such that

\[
f \circ c = \mathrm{D} b = (- 1) ^ {n + 1} b \circ d _ {n + 1}.
\]

We thus have a commutative diagram

\[
\begin{array}{c} L _ {n + 1} (N) \xrightarrow {d _ {n + 1}} L _ {n} (N) \\ (- 1) ^ {n + 1} c \Big \downarrow \\ 0 \longrightarrow M ^ {\prime} \xrightarrow {f} M \xrightarrow {g} M ^ {\prime \prime} \longrightarrow 0 \\ 1 _ {M ^ {\prime}} \Big \downarrow \\ 0 \longrightarrow M ^ {\prime} \xrightarrow {e _ {M ^ {\prime}}} I ^ {0} (M ^ {\prime}) \xrightarrow {\delta^ {0}} I ^ {1} (M ^ {\prime}). \end{array}
\]

But, in \(\operatorname{Homgr}_{\mathbf{A}}(\mathbf{L}(\mathbf{N}), \mathbf{l}(\mathbf{M}'))\) , we have

\[
\mathrm{D} \left(v ^ {0} \circ b\right) = \delta^ {0} \circ v ^ {0} \circ b - (- 1) ^ {n} v ^ {0} \circ b \circ d _ {n + 1} = v ^ {1} \circ a ^ {n} + e _ {\mathrm{M} ^ {\prime}} \circ c.
\]

The classes of \(e_{\mathbf{M}'} \circ c\) and \((-v^1) \circ a\) in \(\operatorname{Ext}_{\mathbf{A}}^{n+1}(\mathbf{N}, \mathbf{M}')\) are equal, whence \(a\) ).

\begin{enumerate}
\def\labelenumi{\alph{enumi})}
\setcounter{enumi}{1}
\tightlist
\item
  Consider a commutative diagram
\end{enumerate}

\[
\begin{array}{c} \mathrm{L} _ {1} (\mathrm{M} ^ {\prime \prime}) \xrightarrow {d _ {1}} \mathrm{L} _ {0} (\mathrm{M} ^ {\prime \prime}) \xrightarrow {p _ {\mathrm{M} ^ {\prime \prime}}} \mathrm{M} ^ {\prime \prime} \longrightarrow 0 \\ \Biggl \downarrow u _ {1} \\ 0 \longrightarrow \mathrm{M} ^ {\prime} \xrightarrow {f} \mathrm{M} \xrightarrow {g} \mathrm{M} ^ {\prime \prime} \longrightarrow 0. \end{array}
\]

By definition, \(\theta\) is the class of \(-u_{1}\in\mathrm{Homgr}_{\mathbf{A}}^{1}(\mathrm{L}(\mathbf{M}^{\prime\prime}),\mathbf{M}^{\prime})\) . Let moreover \(\alpha\in\mathrm{Ext}^{n}(\mathbf{M}^{\prime},\mathbf{N})\) represented by an element \(a\) of \(\mathrm{Homgr}_{\mathbf{A}}^{n}(\mathbf{M}^{\prime},\mathrm{I}(\mathbf{N}))\) . By construction, \(\delta^{n}(\alpha)\) is obtained as follows: one extends \(a^{n}\in\mathrm{Hom}_{\mathbf{A}}(\mathbf{M}^{\prime},\mathrm{I}^{n}(\mathbf{N}))\) to

\[
b \in \operatorname{Hom} _ {\mathbf {A}} (\mathbf {M}, \mathrm{I} ^ {n} (\mathbf {N}))
\]

and \(\delta^n (\alpha)\) is the class of \(c\circ p_{\mathbf{M}''}\) , where \(c\in \mathrm{Hom}_{\mathbf{A}}(\mathbf{M}^{\prime \prime},\mathrm{I}^{n + 1}(\mathbf{N}))\) is such that

\[
g \circ c = D b = \delta^ {n + 1} \circ b.
\]

We thus have a commutative diagram

\[
\begin{array}{c} L _ {1} (M ^ {\prime \prime}) \xrightarrow {d _ {1}} L _ {0} (M ^ {\prime \prime}) \xrightarrow {p _ {M ^ {\prime \prime}}} M ^ {\prime \prime}. \longrightarrow 0 \\ u _ {1} \Biggl \downarrow \quad \quad \quad \quad \quad \quad \quad \quad \quad \quad \quad \quad \quad \quad \quad \quad \quad \quad \quad \quad \quad \quad \quad \quad \quad \quad \quad \quad \quad \quad \quad \quad \quad \quad \quad \quad \quad \quad \quad \quad \quad \quad \quad \quad \quad \quad \quad 1 _ {M ^ {\prime}} \\ 0 \longrightarrow M ^ {\prime} \xrightarrow {f} M \xrightarrow {g} M ^ {\prime \prime} \longrightarrow 0 \\ a ^ {n} \Biggl \downarrow \quad \quad \quad \quad \quad \quad \quad \quad \quad \quad \quad \quad c \Biggl \downarrow \\ I ^ {n} (N) \xrightarrow {\delta^ {n}} I ^ {n + 1} (N). \end{array}
\]

But, in \(\operatorname{Homgr}_{\mathbf{A}}(\mathbf{L}(\mathbf{M}^{\prime \prime}),\mathbf{l}(\mathbf{N}))\) , we have

\[
\mathrm{D} (b \circ u _ {0}) = \delta^ {n} \circ b \circ u _ {0} - (- 1) ^ {n} b \circ u _ {0} \circ c = c \circ p - (- 1) ^ {n} a ^ {n} \circ u _ {1}.
\]

The classes of \(c \circ p\) and of \((-1)^{n+1} a^n \circ (-u_1)\) in \(\operatorname{Ext}_{\mathbf{A}}^{n+1} (\mathbf{M}''\) , \(\mathbf{N}\) ) are therefore equal, whence \(b\) ).

COROLLARY 1. --- a) The connecting homomorphism Hom \(_{A}\) (M'', M'') → Ext \(_{A}^{1}\) (M'', M') sends 1 \(_{M''}\) to θ.

\begin{enumerate}
\def\labelenumi{\alph{enumi})}
\setcounter{enumi}{1}
\tightlist
\item
  The connecting homomorphism \(\operatorname{Hom}_{\mathbf{A}}(\mathbf{M}^{\prime},\mathbf{M}^{\prime})\to\operatorname{Ext}_{\mathbf{A}}^{1}(\mathbf{M}^{\prime\prime},\mathbf{M}^{\prime})\) sends \(1_{M'}\) to \(-\theta\).
\end{enumerate}

COROLLARY 2. --- Consider two exact sequences of left A-modules

\[
0 \rightarrow \mathbf {M} ^ {\prime} \rightarrow \mathbf {M} \rightarrow \mathbf {M} ^ {\prime \prime} \rightarrow 0
\]

\[
0 \rightarrow \mathrm{N} ^ {\prime} \rightarrow \mathrm{N} \rightarrow \mathrm{N} ^ {\prime \prime} \rightarrow 0.
\]

Then the composed homomorphisms of connecting homomorphisms

\[
\operatorname{Ext} _ {\mathbf {A}} ^ {n} \left(\mathrm{M} ^ {\prime}, \mathrm{N} ^ {\prime \prime}\right)\rightarrow \operatorname{Ext} _ {\mathbf {A}} ^ {n + 1} \left(\mathrm{M} ^ {\prime \prime}, \mathrm{N} ^ {\prime \prime}\right)\rightarrow \operatorname{Ext} _ {\mathbf {A}} ^ {n + 2} \left(\mathrm{M} ^ {\prime \prime}, \mathrm{N} ^ {\prime}\right)
\]

and

\[
\operatorname{Ext} _ {\mathbf {A}} ^ {n} \left(\mathrm{M} ^ {\prime}, \mathrm{N} ^ {\prime \prime}\right)\rightarrow \operatorname{Ext} _ {\mathbf {A}} ^ {n + 1} \left(\mathrm{M} ^ {\prime}, \mathrm{N} ^ {\prime}\right)\rightarrow \operatorname{Ext} _ {\mathbf {A}} ^ {n + 2} \left(\mathrm{M} ^ {\prime \prime}, \mathrm{N} ^ {\prime}\right)
\]

are opposite.

Indeed, if \(\theta_{1}\) , \(\theta_{2}\) are the classes associated with the given exact sequences, and if \(\alpha \in \operatorname{Ext}_{\mathbf{A}}^{n}(\mathbf{M}', \mathbf{M}'')\) the images of \(\alpha\) are respectively

\[
\theta_ {2} \circ ((- 1) ^ {n + 1} \alpha \circ \theta_ {1}) \quad \text { and } \quad (\theta_ {2} \circ \alpha) \circ ((- 1) ^ {n + 2} \theta_ {1}).
\]

Consider an exact sequence of left A-modules

\[
0 \rightarrow \mathrm{N} \rightarrow \mathrm{R} _ {n} \xrightarrow {f _ {n}} \mathrm{R} _ {n - 1} \xrightarrow {f _ {n - 1}} \dots \rightarrow \mathrm{R} _ {1} \xrightarrow {f _ {1}} \mathrm{M} \rightarrow 0\tag{S}
\]

and set \(K_{0} = M\) , \(K_{i} = Ker f_{i}, i = 1, \ldots, n - 1\) , \(K_{n} = N\) . We therefore have exact sequences

\[
0 \to \mathbf {K} _ {i} \to \mathbf {R} _ {i} \to \mathbf {K} _ {i - 1} \to 0, \quad 1 \leqslant i \leqslant n,\tag{9}
\]

to which are associated, for every left A-module P, connecting homomorphisms

\[
\begin{array}{r l}&{\mathrm{Ext} _ {\mathbf {A}} ^ {m} \left(\mathrm{P}, \mathbf {K} _ {i - 1}\right)\rightarrow \mathrm{Ext} _ {\mathbf {A}} ^ {m + 1} \left(\mathrm{P}, \mathbf {K} _ {i}\right),}\\&{\qquad \mathrm{Ext} _ {\mathbf {A}} ^ {m} \left(\mathbf {K} _ {i}, \mathrm{P}\right)\rightarrow \mathrm{Ext} _ {\mathbf {A}} ^ {m + 1} \left(\mathbf {K} _ {i - 1}, \mathrm{P}\right),}\end{array}
\]

whence by composition of the iterated connecting homomorphisms, associated with (S)

\[
\begin{array}{l} \delta^ {m} (\mathrm{P}, \mathcal {S}): \operatorname{Ext} _ {\mathrm{A}} ^ {m} (\mathrm{P}, \mathrm{M}) \to \operatorname{Ext} _ {\mathrm{A}} ^ {m + n} (\mathrm{P}, \mathrm{N}) \\ \delta^ {m} (\mathcal {S}, \mathrm{P}): \operatorname{Ext} _ {\mathrm{A}} ^ {m} (\mathrm{N}, \mathrm{P}) \to \operatorname{Ext} _ {\mathrm{A}} ^ {m + n} (\mathrm{M}, \mathrm{P}). \end{array}
\]

COROLLARY 3. --- If \(\theta\in\mathrm{Ext}_{\mathbf{A}}^{n}(\mathbf{M},\mathbf{N})\) is the class of the exact sequence \((\mathcal{S})\) , we have

\[
\delta^ {m} (\mathrm{P}, \mathscr {S}) (\alpha) = \theta \circ \alpha , \quad \delta^ {m} (\mathscr {S}, \mathrm{P}) (\beta) = (- 1) ^ {m n + n (n + 1) / 2} \beta \circ \theta .
\]

If \(\theta_{i}\in\mathrm{Ext}_{\mathbf{A}}^{1}(\mathbf{K}_{i-1},\mathbf{K}_{i})\) is the class associated with the exact sequence (9), we have by prop. 5

\[
\begin{array}{l} \delta^ {m} (\mathrm{P}, \mathcal {S}) (\alpha) = \theta_ {n} \circ \dots \circ \theta_ {2} \circ \theta_ {1} \circ \alpha \\ \delta^ {m} (\mathcal {S}, \mathrm{P}) (\beta) = (- 1) ^ {(m + 1) + \dots + (m + n)} \beta \circ \theta_ {n} \circ \dots \circ \theta_ {1}. \end{array}
\]

Moreover, by prop. 3 (X, p.~118), we have \(\theta = \theta_{n} \circ \cdots \circ \theta_{1}\) . The corollary follows immediately from this, and from the relation (E, III, p.~44)

\[
(m + 1) + \dots + (m + n) = m n + n (n + 1) / 2.
\]

COROLLARY 4. --- If each module \(\mathbf{R}_i\) , \(i = 1, \ldots, n\) , is injective (resp. projective), the map \(\alpha \mapsto \theta \circ \alpha\) (resp. \(\alpha \mapsto \alpha \circ \theta\) ) from \(\operatorname{Ext}_{\mathbf{A}}^{m}(\mathbf{P}, \mathbf{M})\) to \(\operatorname{Ext}_{\mathbf{A}}^{m+n}(\mathbf{P}, \mathbf{N})\) (resp. from \(\operatorname{Ext}_{\mathbf{A}}^{m}(\mathbf{N}, \mathbf{P})\) to \(\operatorname{Ext}_{\mathbf{A}}^{m+n}(\mathbf{M}, \mathbf{P})\) ) is bijective for every A-module P and every integer \(m > 0\) . This follows indeed from cor. 3 and the exact sequences

\[
\operatorname{Ext} _ {\mathbf {A}} ^ {m + i - 1} \left(\mathrm{P}, \mathrm{R} _ {i}\right)\rightarrow \operatorname{Ext} _ {\mathbf {A}} ^ {m + i - 1} \left(\mathrm{P}, \mathrm{K} _ {i - 1}\right)\rightarrow \operatorname{Ext} _ {\mathbf {A}} ^ {m + i} \left(\mathrm{P}, \mathrm{K} _ {i}\right)\rightarrow \operatorname{Ext} _ {\mathbf {A}} ^ {m + i} \left(\mathrm{P}, \mathrm{R} _ {i}\right)
\]

\[
\left(\text { resp. } \operatorname{Ext} _ {\mathbf {A}} ^ {m + i - 1} \left(\mathrm{R} _ {i}, \mathrm{P}\right)\rightarrow \operatorname{Ext} _ {\mathbf {A}} ^ {m + i - 1} \left(\mathrm{K} _ {i}, \mathrm{P}\right)\rightarrow \operatorname{Ext} _ {\mathbf {A}} ^ {m + i} \left(\mathrm{K} _ {i - 1}, \mathrm{P}\right)\rightarrow \operatorname{Ext} _ {\mathbf {A}} ^ {m + i} \left(\mathrm{R} _ {i}, \mathrm{P}\right)\right),
\]

whose extreme terms are zero by hypothesis.

Remark. --- The definitions and propositions of nos. 3 to 6 apply to right A-modules, considered as left modules over the opposite ring A° of A.

